\documentclass[11pt]{article}
\usepackage[a4paper,margin=2.8cm]{geometry}
\usepackage{microtype}
\usepackage{parskip}
\begin{document}

\noindent
[Date]

\noindent
The Editors\\
\emph{Nature Machine Intelligence}

\bigskip
\noindent
Dear Editors,

We are pleased to submit our manuscript, \textbf{``Useful possibility
collapse: an intervention-aware framework for adaptive emergence in
learning systems''}, for consideration as an Article in \emph{Nature
Machine Intelligence}.

``Emergence'' is measured differently in complex-systems science
(causal efficacy of macro-variables), information theory (synergy),
and machine learning (abrupt ability curves)---and the communities
disagree about what counts. The dispute is not academic: whether large
models have ``emergent abilities'' is currently argued by reading
performance curves, although a curve alone carries no information
about where an ability came from.

Our manuscript addresses this gap in three steps.

\textbf{First, an honest theory architecture.} The trajectory-space
quantity we start from is a known specific-information object (the
Bayesian-surprise divergence), and we position it as such: it fixes
what emergence claims must be about---distributions over futures---but
does not by itself carry a verdict. The verdict comes from two layered
criteria stated on a declared observer contract: \emph{structural
collapse} (potential, conditional selectivity, specificity) and
\emph{adaptive qualification} (usefulness, endogeneity, acquisition).
The layering is testable and tested: under a second plausible observer
contract, all 60 control verdicts and 14/15 structural verdicts are
unchanged, while value verdicts flip conservatively in five cases,
exactly the contract-relativity the definition declares. The signatures
audited here---performance jumps, representation jumps, PID-style
synergy, Rosas' $\Psi$, and Hoel's effective information---are lossy
projections of the same trajectory-law family; we compute the latter
two from their published equations with zero Monte-Carlo error on
fully enumerated Markov chains, within declared candidate families,
and we do not claim these formalisms are invalid in their home use
cases.

\textbf{Second, comparisons with matched degrees of freedom.} Beyond
single-signal audits, we give the prior signatures equal or greater
freedom---multi-threshold conjunctions, logistic regression and
decision trees over all signals including the exact EI and
$\Psi$---and test frozen transfer to fresh systems. Every fitted
baseline caps at 0.9 and misclassifies the true conditional emergence,
while the frozen six-component protocol scores 10/10 on the same fresh
seed. Two independent domains also exposed the same methodological
principle, which we believe will matter for capability auditing: in
converged systems, the behaving policy prices its own triggers into
the baseline, so useful emergence is visible through counterfactual
contrast rather than observed improvement along the behaving
trajectory.

\textbf{Third, stress tests across complementary evidence types.} Under
internally frozen author-maintained protocols (with all thirty-one
recorded misses retained), the full six-component protocol is tested in
three built domains and one public benchmark. It classifies 25/25 fresh
external-swarm verdicts; in Contextual Level-Based Foraging, 9/10
learned policies pass while all 40 initialization/scripted controls
reject; and in a controlled latent-context sequence task, 10/10
next-token-trained transformers pass while all 40 controls reject.
Decisively, in \emph{unmodified} Overcooked-AI---a public third-party
coordination benchmark---the complete protocol, thresholds and five
predictions were frozen with a \emph{public external timestamp} that
predates every confirmatory measurement, and all five registered
predictions passed: learned self-play policies are accepted on exactly
the registered 8/12 seed line, all 48 control verdicts are rejections
with the failure routes the layered definition predicts, and the
primary seed-level usefulness do-contrast is positive on 12/12 seeds
($p=2\times10^{-4}$). The retained misses are used to delimit
scope rather than repaired. For example, persistence is stable under
declared temporal and observational perturbations but not under novel
geometries, and the sequence model generalizes to distractors but not to
unseen marker positions. Separately, a frozen future-distribution score
prospectively discovers value-critical decisions in uncurated chess
positions across two months (AUROC 0.730 and 0.725), transfers to a
classical-versus-NNUE referee, and is directionally consistent under an
engine-free realized-outcome referee whose registered interaction we
report as an underpowered null. The process proxy recovers grokking and
induction-head transitions and transfers through Pythia-1.4B, while its
registered 2.8B and thinning failures expose checkpoint-resolution limits.

We believe the work fits \emph{Nature Machine Intelligence}'s scope
directly: it turns ``did an ability emerge?'' from a curve-reading
exercise into an intervention-aware measurement programme, with relevance
to capability auditing and forecasting in machine learning, and with a
transparency standard (frozen thresholds, registered failures, full
code and outputs) that we hope raises the bar for definitional claims
in this area.

The manuscript is not under consideration elsewhere. All code,
outputs, prospectively frozen protocols, a reproducibility manifest,
invalid-data registry, and a complete prediction ledger sufficient to
audit every figure accompany the submission.

Suggested referees: [names spanning the three communities---causal
emergence, information decomposition, and LLM capability
evaluation---to be filled in].

\bigskip
\noindent
Sincerely,

\noindent
[Corresponding author]\\
on behalf of all authors

\end{document}
